\documentclass[11pt]{article}

\usepackage[margin=1in]{geometry}
\usepackage[T1]{fontenc}
\usepackage[utf8]{inputenc}
\usepackage{lmodern}
\usepackage{microtype}
\usepackage{amsmath}
\usepackage{hyperref}
\usepackage{enumitem}
\usepackage{xcolor}

\hypersetup{
  colorlinks=true,
  linkcolor=blue!50!black,
  urlcolor=blue!50!black
}

\title{\textbf{Future Work Note}\\[4pt]
Sequential Pose-Aware Localization with Kalman Filtering}
\author{}
\date{\today}

\begin{document}
\maketitle

\section*{Core Idea}
The idea being suggested is a shift from \textbf{frame-by-frame place recognition} to \textbf{state estimation over time}.

Right now, the VPR pipeline mostly asks:
\begin{quote}
``Given this current image, which reference image looks most similar?''
\end{quote}

The proposed future direction is closer to:
\begin{quote}
``Given where the robot was a moment ago, how it moved, and what it is seeing now, where should it be in the map?''
\end{quote}

This is where a Kalman-filter-style idea comes in.

\section*{Intuition}
A robot should not localize from vision alone at every frame as if each frame were independent. In navigation, the robot has temporal continuity:

\begin{itemize}[leftmargin=*]
  \item it has a previous estimated pose,
  \item it has some notion of motion between frames,
  \item it knows that the map was built in a 2D world,
  \item so it can predict what location and orientation it is likely to have next.
\end{itemize}

Instead of matching the current frame against the entire map blindly, the system first forms an \textbf{expectation} of where the robot should be, then uses the visual observation to correct that expectation.

\section*{State Representation}
The hidden state would be the robot pose, for example:
\[
(x, y, \theta)
\]
where:
\begin{itemize}[leftmargin=*]
  \item $x$ is the robot position in the map,
  \item $y$ is the robot position in the map,
  \item $\theta$ is the robot orientation.
\end{itemize}

This means the system is not only estimating ``which image matches'' but also:
\begin{quote}
``The robot is probably here, facing this direction.''
\end{quote}

\section*{Prediction and Update}
Conceptually, a Kalman filter works in two phases.

\subsection*{1. Prediction}
Use the previous pose estimate and the robot's motion to predict the new pose.

For example, if the robot moved forward slightly and rotated a little, the system predicts a new $(x, y, \theta)$.

\subsection*{2. Update}
Use the current visual observation to correct that prediction.

For example, if the VPR system says the current view looks most similar to reference images around a particular campus location, the pose estimate is pulled toward that region.

So the filter combines:
\begin{itemize}[leftmargin=*]
  \item \textbf{motion expectation}, and
  \item \textbf{visual evidence}.
\end{itemize}

This is much stronger than using visual similarity alone.

\section*{What ``Expectation'' Means Here}
The professor's phrase about comparing \textbf{what the robot is seeing} versus \textbf{what it expects to see} is important.

The expectation comes from:
\begin{itemize}[leftmargin=*]
  \item the estimated pose,
  \item the motion model,
  \item the map.
\end{itemize}

If the robot believes it is near a certain place and facing a certain direction, then it should expect visual evidence consistent with nearby reference views, not arbitrary places in the database.

So, instead of searching the entire map equally, the system could:
\begin{itemize}[leftmargin=*]
  \item restrict matching to nearby candidates,
  \item weight nearby candidates more strongly,
  \item compare the observation to what should be visible from the predicted pose,
  \item smooth decisions over a sequence rather than treating each frame independently.
\end{itemize}

\section*{Why This Helps VPR}
This matters because place-recognition errors often happen when:
\begin{itemize}[leftmargin=*]
  \item two places look visually similar,
  \item viewpoint changes affect the descriptor,
  \item lighting changes make matching noisy,
  \item a single frame is ambiguous.
\end{itemize}

A sequential filter helps because it introduces a physical constraint:
\begin{quote}
``Even if one frame is confusing, the robot cannot teleport. Its next likely pose must be near its previous pose.''
\end{quote}

That constraint reduces false matches and makes the system more navigation-oriented.

\section*{How This Fits Our Project}
A natural future-work version for this project would be:

\begin{enumerate}[leftmargin=*]
  \item During initial mapping, assign each reference image an approximate 2D pose or sequence index.
  \item During live inference, maintain an estimated robot pose.
  \item Use robot motion or frame-to-frame progression to predict the next pose.
  \item Use VPR outputs as the observation model.
  \item Fuse motion and visual evidence with a Kalman filter or a related Bayesian filter.
\end{enumerate}

Then the system would no longer only say:
\begin{quote}
``The best visual match is reference image 42.''
\end{quote}

It would instead say something closer to:
\begin{quote}
``Based on motion, the robot should be near pose $(x, y, \theta)$, and the current visual observation supports reference images around that pose.''
\end{quote}

That is much closer to a real localization system for navigation.

\section*{Important Technical Note}
Strictly speaking, a basic linear Kalman filter may not be sufficient for the full problem, because:
\begin{itemize}[leftmargin=*]
  \item robot motion and orientation are nonlinear,
  \item visual observations can be ambiguous or multimodal.
\end{itemize}

So in a more complete system, related methods might be more appropriate, such as:
\begin{itemize}[leftmargin=*]
  \item Extended Kalman Filter (EKF),
  \item Unscented Kalman Filter (UKF),
  \item particle filters,
  \item other Bayesian filtering or pose-graph approaches.
\end{itemize}

However, as a \textbf{future work} direction, ``Kalman-filter-based fusion of motion prediction and VPR observations'' is the correct high-level idea and a strong way to frame the concept.

\section*{Polished Future Work Paragraph}
As future work, we plan to move from frame-wise visual place recognition to \textbf{sequential pose-aware localization}. Instead of treating each query image independently, the system would maintain an estimate of the robot's 2D pose and orientation over time. A motion model would predict the robot's next state, while VPR outputs would act as visual observations used to correct that prediction. This would allow the robot to compare what it currently sees against what it is expected to see from its predicted location, improving robustness to viewpoint changes, perceptual ambiguity, and noisy single-frame matches. A Kalman-filter-based formulation, or a related nonlinear Bayesian filtering approach, could provide a principled way to fuse motion and visual evidence for navigation-oriented localization.

\end{document}
